\subsection{Modele i pokrycie słownika}\label{sec:we-models}

Osadzenia policzono dla wszystkich \weNWords{} słów (\weNTokens{} wystąpień; bez interpunkcji i samych liczb)
z \weNSentences{} zaanotowanych zdań. Porównano dwa modele word2vec~\cite{word2vec} (skip-gram, 100 wymiarów),
które różnią się korpusem uczącym:
\begin{itemize}
  \item \textbf{IPI PAN} -- gotowy model \texttt{nkjp+wiki-forms-all-100-skipg-ns}~\cite{dsmodels}, uczony na pełnym
        Narodowym Korpusie Języka Polskiego i polskiej Wikipedii; wczytano \weVocabIpi{} najczęstszych form;
  \item \textbf{PolEmo} -- model wytrenowany przez nas w bibliotece gensim na wszystkich \wePolemoSentences{} zdaniach
        PolEmo (\wePolemoWords{} słów; słowa występujące co najmniej 3 razy, 15 epok); słownik \weVocabPol{} słów.
\end{itemize}
Wektor słowa wyszukiwano kolejno w formie małymi literami, w oryginalnej pisowni i z wielką literą na początku.

\begin{table}[H]
\centering\small
\resizebox{\textwidth}{!}{\begin{tabular}{lrrrrrrrr}
\toprule
Model & słowa & wystąpienia & słowa zaanotowane & WYR\_POZ ($n$=145) & WYR\_NEG ($n$=86) & WZM ($n$=19) & OSL ($n$=13) & ODW ($n$=4) \\
\midrule
IPI PAN & 94{,}6\% & 97{,}1\% & 94{,}9\% & 94\% & 91\% & 100\% & 100\% & 100\% \\
PolEmo & 85{,}2\% & 92{,}0\% & 83{,}7\% & 86\% & 65\% & 95\% & 100\% & 100\% \\
\bottomrule
\end{tabular}}
\caption{Pokrycie słownika: odsetek słów (wystąpień) zaanotowanego korpusu, dla których model ma wektor.
„Słowa zaanotowane'' -- z co najmniej jedną anotacją różną od O; kolumny etykiet -- słowa o danej etykiecie
(najczęstsza etykieta anotacji, jeśli co najmniej połowa anotacji słowa jest różna od O).}
\label{tab:we-coverage}
\end{table}

Brakujące wektory w modelu IPI PAN to głównie zapisy bez polskich znaków (\emph{lazience} zamiast \emph{łazience}), formy z cyframi (\emph{2osobowym}), słowa rzadkie (\emph{acellularne}).
